\chapter{Kommentierte Quellen}

Achtundzwanzig Werke, eingeordnet: Zu jedem steht, wofür man es aufschlägt und was es leistet, das andere nicht leisten.

\lit{1}{Agresti, A.: \textit{Categorical Data Analysis}. 3.\ Aufl., Wiley 2013.}{Das Standardwerk zu allem, was in Kapitel 11 nur angetippt wurde: logistische Modelle, ordinale Regression, Verfahren für Vierfeldertafeln. Die Abschnitte über exakte Verfahren bei kleinen Zahlen sind besonders sorgfältig.}

\lit{2}{Altman, D.\,G.: \textit{Practical Statistics for Medical Research}. Chapman \& Hall 1991.}{\textbf{Das Buch für die medizinische Anwendung.} Altman schreibt für Ärztinnen und Ärzte, nicht für Statistiker, und er schreibt ausgezeichnet. Trotz des Alters kaum überholt --- die Fallen, vor denen er warnt, sind dieselben geblieben. Altman war einer der Urheber der CONSORT-Empfehlungen.}

\lit{3}{Amrhein, V., Greenland, S., McShane, B.: \textit{Retire statistical significance}. Nature 567, 2019.}{Zwei Seiten, über achthundert Unterzeichner. Der Aufruf, das Wort \enquote{signifikant} aufzugeben --- mit dem Hauptargument, dass die Dichotomie selbst der Fehler ist, gleichgültig wo die Schwelle liegt. Frei verfügbar, in einer halben Stunde gelesen.}

\lit{4}{Benjamini, Y., Hochberg, Y.: \textit{Controlling the False Discovery Rate}. J.\ Royal Stat.\ Soc.\ B 57, 1995.}{Die Arbeit hinter Kapitel 14. Sie führt einen neuen Fehlerbegriff ein und zeigt, dass er bei vielen Tests der brauchbarere ist. Der Aufsatz war zunächst wenig beachtet und wurde mit der Genomforschung zum meistzitierten Verfahren seiner Art.}

\lit{5}{Bland, M.: \textit{An Introduction to Medical Statistics}. 4.\ Aufl., Oxford 2015.}{Die freundlichste Einführung des Gebiets, mit vielen durchgerechneten medizinischen Beispielen und einem ausführlichen Aufgabenteil. Wo dieser Band zu knapp war, ist Bland der nächste Schritt.}

\lit{6}{Borenstein, M., Hedges, L., Higgins, J., Rothstein, H.: \textit{Introduction to Meta-Analysis}. 2.\ Aufl., Wiley 2021.}{\textbf{Für Kapitel 14 die maßgebliche Darstellung.} Behandelt festes und zufälliges Modell, Heterogenität, Untergruppen und Metaregression --- durchgehend mit dem ausdrücklichen Hinweis, welches Modell welche Frage beantwortet.}

\lit{7}{Casella, G., Berger, R.: \textit{Statistical Inference}. 2.\ Aufl., Duxbury 2002.}{Die theoretische Grundlage zu allem, was dieser Band anwendet: Verteilungstheorie, Schätztheorie, Testtheorie, vollständig bewiesen. Anspruchsvoll; hier stehen auch die Herleitungen, die dieser Band ausgelassen hat.}

\lit{8}{Higgins, J. u.\,a.\ (Hrsg.): \textit{Cochrane Handbook for Systematic Reviews of Interventions}. Version 6.5, 2024.}{Die Arbeitsanleitung der Cochrane-Zusammenarbeit --- frei im Netz verfügbar und fortlaufend aktualisiert. Kein Lehrbuch, sondern ein Handbuch: Wie sucht man Studien, wie bewertet man ihr Verzerrungsrisiko, wie fasst man zusammen.}

\lit{9}{Cox, D.\,R.: \textit{Regression Models and Life-Tables}. J.\ Royal Stat.\ Soc.\ B 34, 1972.}{Die meistzitierte Arbeit der Statistik überhaupt. Lesenswert auch wegen der abgedruckten Diskussion, in der mehrere angesehene Statistiker bezweifeln, dass das vorgeschlagene Verfahren überhaupt zulässig sei.}

\lit{10}{Efron, B., Hastie, T.: \textit{Computer Age Statistical Inference}. Cambridge 2016.}{\textbf{Das beste Buch über den Zusammenhang von klassischer Statistik und maschinellem Lernen.} Efron ordnet beide Traditionen historisch ein und zeigt, wie Rechenleistung die Theorie verändert hat. Frei verfügbar.}

\lit{11}{Efron, B., Tibshirani, R.: \textit{An Introduction to the Bootstrap}. Chapman \& Hall 1993.}{Die Darstellung des Verfahrens durch seinen Urheber, ungewöhnlich zugänglich geschrieben. Enthält auch die verbesserten Intervallverfahren, die dieser Band nur erwähnt hat.}

\lit{12}{Fisher, R.\,A.: \textit{The Design of Experiments}. Oliver \& Boyd 1935.}{Das Buch, das Randomisierung und Signifikanztest begründete. Das erste Kapitel enthält den Teeversuch aus Kapitel 8. Lesenswert als Quelle und als Beispiel dafür, wie unbekümmert Fisher über Fragen hinwegging, die später zu Streitpunkten wurden.}

\lit{13}{Gelman, A., Carlin, J., Stern, H. u.\,a.: \textit{Bayesian Data Analysis}. 3.\ Aufl., Chapman \& Hall 2013.}{Das Standardwerk des bayesschen Zugangs, der in diesem Band nur angedeutet wurde. Der Zugang ist ungewöhnlich praktisch: Es geht weniger um die Grundlagenfrage als um das Rechnen mit wirklichen Daten. Frei verfügbar.}

\lit{14}{Gigerenzer, G.: \textit{Das Einmaleins der Skepsis}. Berliner Taschenbuch Verlag 2004.}{\textbf{Auf Deutsch, für alle lesbar, und in der Sache streng.} Gigerenzer zeigt, wie sich Wahrscheinlichkeiten verständlich darstellen lassen --- nämlich als natürliche Häufigkeiten, genau wie in Kapitel 13. Enthält die Untersuchungen, wonach die Mehrheit der befragten Ärzte den positiven prädiktiven Wert falsch einschätzt.}

\lit{15}{Good, P.: \textit{Permutation, Parametric and Bootstrap Tests of Hypotheses}. 3.\ Aufl., Springer 2005.}{Für Kapitel 9 die Vertiefung. Behandelt Permutationstests systematisch, einschließlich der Fälle, in denen sie versagen --- und der Frage, wie man daraus Konfidenzintervalle gewinnt.}

\lit{16}{Greenland, S. u.\,a.: \textit{Statistical tests, P values, confidence intervals, and power: a guide to misinterpretations}. Eur.\ J.\ Epidemiol.\ 31, 2016.}{\textbf{Fünfundzwanzig Fehldeutungen, einzeln widerlegt.} Die ausführliche Fassung der Warnkästen aus Kapitel 6 und 7, von den maßgeblichen Fachleuten geschrieben. Frei verfügbar und für jeden lesbar, der diesen Band durchgearbeitet hat.}

\lit{17}{Harrell, F.: \textit{Regression Modeling Strategies}. 2.\ Aufl., Springer 2015.}{Die Antwort auf die Frage, was man statt schrittweiser Variablenauswahl tun soll. Harrell begründet ausführlich, warum die verbreiteten Vorgehensweisen schaden, und stellt Alternativen vor. Sein Begleitblog ist eine eigene Fundgrube.}

\lit{18}{Hernán, M., Robins, J.: \textit{Causal Inference: What If}. Chapman \& Hall 2020.}{\textbf{Das Buch zur kausalen Inferenz --- frei verfügbar und außerordentlich klar.} Entwickelt den Begriff der Zielstudie aus Kapitel 15 und zeigt an Beispielen, wie sich Widersprüche zwischen Beobachtungs- und randomisierten Studien auflösen lassen. Der erste Teil verlangt kaum Vorkenntnisse.}

\lit{19}{Hosmer, D., Lemeshow, S., Sturdivant, R.: \textit{Applied Logistic Regression}. 3.\ Aufl., Wiley 2013.}{Die anwendungsnahe Darstellung zu Kapitel 11, mit ausführlicher Behandlung der Modellprüfung --- des Teils, der in der Praxis am häufigsten unterbleibt.}

\lit{20}{Ioannidis, J.: \textit{Why Most Published Research Findings Are False}. PLoS Medicine 2, 2005.}{Sechs Seiten, die ein Fach erschüttert haben. Die Rechnung ist elementar --- es ist die Vierfeldertafel aus Aufgabe 7.3 --- und ihre Folgerungen sind unbequem. Frei verfügbar.}

\lit{21}{Kaplan, E., Meier, P.: \textit{Nonparametric Estimation from Incomplete Observations}. J.\ Amer.\ Stat.\ Assoc.\ 53, 1958.}{Die Arbeit hinter Kapitel 11, entstanden aus zwei unabhängigen Einreichungen und vier Jahren Verhandlung. Sie ist erstaunlich gut lesbar und begründet ausführlich, warum die naheliegenden Alternativen scheitern.}

\lit{22}{Klein, J., Moeschberger, M.: \textit{Survival Analysis}. 2.\ Aufl., Springer 2003.}{Die vollständige Darstellung der Überlebenszeitanalyse: Zensierungsarten, Log-Rang-Varianten, Cox-Modell mit zeitabhängigen Größen, konkurrierende Ereignisse.}

\lit{23}{Pearl, J., Mackenzie, D.: \textit{The Book of Why}. Basic Books 2018.}{Die allgemeinverständliche Darstellung dessen, was in Kapitel 3 und 15 nur angedeutet wurde. Deutsch als \kursiv{Wenn Sie mich fragen} erschienen. Pearl argumentiert streitlustig und gelegentlich einseitig --- das macht das Buch lesbar und verlangt Mitdenken.}

\lit{24}{Schumacher, M., Schulgen, G.: \textit{Methodik klinischer Studien}. 3.\ Aufl., Springer 2008.}{\textbf{Auf Deutsch, und für die Planung von Studien die praktischste Quelle.} Behandelt Studientypen, Fallzahlplanung, Zwischenauswertungen und Auswertungspläne --- also Kapitel 12 und 14 in der Ausführlichkeit, die man braucht, wenn man selbst plant.}

\lit{25}{Senn, S.: \textit{Statistical Issues in Drug Development}. 3.\ Aufl., Wiley 2021.}{Senn schreibt scharf, kenntnisreich und gegen verbreitete Auffassungen. Die Kapitel über Baseline-Vergleiche, Untergruppenanalysen und Cross-over-Studien räumen mit mehreren hartnäckigen Irrtümern auf.}

\lit{26}{Tukey, J.: \textit{Exploratory Data Analysis}. Addison-Wesley 1977.}{Das Buch, das den Boxplot einführte und die erkundende Auswertung als eigenständige Tätigkeit begründete. Tukeys Haltung --- erst hinsehen, dann rechnen --- ist die Grundlage des Anscombe-Abschnitts in Kapitel 2.}

\lit{27}{Wasserstein, R., Lazar, N.: \textit{The ASA Statement on p-Values}. The American Statistician 70, 2016.}{Die Erklärung aus Kapitel 7, mit ausführlicher Begründung jedes der sechs Punkte und einer lesenswerten Vorgeschichte. Frei verfügbar; zwei Seiten Text und zwanzig Seiten Kommentare.}

\lit{28}{Weiß, C.: \textit{Basiswissen Medizinische Statistik}. 8.\ Aufl., Springer 2024.}{Die verbreitetste deutschsprachige Einführung, auf das Medizinstudium zugeschnitten. Knapper als dieser Band und näher an den Prüfungsanforderungen. Als Nachschlagewerk und zur Wiederholung brauchbar.}

\begin{merksatz}
\boxtitel{Wenn Sie nur drei Bücher kaufen}
\begin{center}
\renewcommand{\arraystretch}{1.6}
\begin{tabular}{@{}P{1.0cm}P{2.8cm}P{4.4cm}@{}}
\toprule
\textbf{[2]} & \textbf{Altman} & \textbf{Für die Anwendung}\\
\textbf{[18]} & \textbf{Hernán, Robins} & Für das Denken in Ursachen\\
\textbf{[10]} & \textbf{Efron, Hastie} & \textbf{Für die Einordnung}\\
\bottomrule
\end{tabular}
\end{center}

\begin{center}
\begin{minipage}{0.88\textwidth}
\kursiv{Altman ist das Buch, das neben dem Schreibtisch liegt --- es beantwortet die Fragen, die beim Auswerten wirklich auftreten. Hernán und Robins ändern die Art, wie man über Daten nachdenkt: Nach den ersten Kapiteln fragt man bei jeder Zahl, welche gedachte Studie sie eigentlich nachbildet. Efron und Hastie ordnen alles ein --- man versteht danach, warum die Statistik so aussieht, wie sie aussieht.}
\end{minipage}
\end{center}

\begin{center}
\begin{minipage}{0.88\textwidth}
\kursiv{Drei Ergänzungen, alle frei verfügbar und alle kurz: [16] Greenland u.\,a.\ mit fünfundzwanzig widerlegten Fehldeutungen, [20] Ioannidis mit sechs Seiten, die man nicht vergisst, und [3] Amrhein u.\,a.\ mit zwei. Auf Deutsch: [14] Gigerenzer und, für die Studienplanung, [24] Schumacher und Schulgen.}
\end{minipage}
\end{center}
\end{merksatz}

\bigskip
\begin{flushright}
\itshape\color{grau}
Die Formel ist das Leichte.\\
Die Frage, ob sie hier gilt,\\
ist das Eigentliche.\\[10pt]
\normalfont\color{grau}\upshape
--- der Anfang dieses Bandes\\
\itshape und sein Maßstab.
\end{flushright}
